\chapter{Solid Modeling: Part and PartDesign}
\label{ch:part-partdesign}

FreeCAD offers two complementary solid-modeling paradigms. Knowing which
to reach for --- and how they differ --- saves a great deal of grief,
especially when the solid is destined for meshing and analysis.

\section{Part: direct/CSG modeling}

The \textsf{Part} workbench exposes OCCT directly. You create
\textbf{primitives} (box, cylinder, sphere, cone, torus) and combine them
with \textbf{Boolean operations} --- union (\emph{fuse}), difference
(\emph{cut}), intersection (\emph{common}). This is
\textbf{constructive solid geometry (CSG)}: the model is a tree of
Booleans over primitives.

Part modeling is fast, robust and scriptable, and it happily produces
multiple independent solids. It is ideal for simple analysis geometry
(a plate with holes, a bracket approximated by boxes and cylinders) and
for programmatic model generation.

\section{PartDesign: feature-based body modeling}

The \textsf{PartDesign} workbench enforces a stricter, more powerful
discipline aimed at manufacturable single parts. Its central object is
the \textbf{Body}, a container that holds an ordered stack of
\textbf{features} producing exactly \emph{one} contiguous solid. Features
are either \emph{additive} (Pad, Revolution, Loft, Sweep) or
\emph{subtractive} (Pocket, Hole, Groove) and are typically driven by
sketches.

Key PartDesign concepts:

\begin{description}
  \item[The tip.] The Body has a \emph{tip} feature marking the current
        end of the feature stack; the tip's solid is the Body's result.
  \item[Sketch-on-face.] Later sketches can be placed on faces created by
        earlier features, chaining design intent.
  \item[Dress-up features.] Fillet, chamfer and draft operate on edges
        of the running solid.
  \item[Datum features.] Datum planes, lines and points provide
        references for sketches and features independent of existing
        geometry.
\end{description}

The single-solid rule is what distinguishes PartDesign from Part: a
PartDesign Body must remain one connected lump at every step, which
guarantees a clean, manufacturable, and \emph{meshable} result.

\section{Choosing a paradigm for analysis geometry}

\begin{center}
\begin{tabularx}{\linewidth}{lXX}
  \toprule
   & \textsf{Part} (CSG) & \textsf{PartDesign} (features) \\
  \midrule
  Best for      & primitives, Booleans, scripted or multi-solid models
                & manufacturable single parts, design-intent editing \\
  Result        & any number of solids & exactly one solid per Body \\
  Editability   & edit primitive params / Boolean tree
                & edit any feature or its sketch, reorder the stack \\
  Analysis note & simplest robust geometry for a quick FEA
                & realistic parts; watch for tiny fillets/faces \\
  \bottomrule
\end{tabularx}
\end{center}

\begin{fcwarn}[Geometry that meshes well]
The mesher (\cref{ch:meshing-convergence}) discretises the B-rep faces of
the solid. Slivers, tiny fillets, and near-degenerate faces force
locally tiny elements or produce poor-quality ones. \emph{Defeature}
before meshing: suppress cosmetic fillets and holes that do not affect
the load path, and avoid unnecessary Boolean fragmentation. A clean
solid is a prerequisite for a trustworthy analysis, not an afterthought.
\end{fcwarn}

\section{From geometry to analysis}

Whichever paradigm you use, the deliverable for Part~III is a single,
clean solid (or a well-defined set of solids with contact interfaces)
whose faces and edges you can select to apply loads and restraints. With
that in hand, we turn to the theory that a finite-element solver brings
to bear on it.
